% ------------------------------------------------------------------------
% file `v3-exercise-2-solution-body.tex'
%   in folder `xsim/'
%
%     solution of type `exercise' with id `2'
%
% generated by the `solution' environment of the
%   `xsim' package v0.20c (2021/02/03)
% from source `v3' on 2022/06/03 on line 52
% ------------------------------------------------------------------------
    1 point par justification correcte.
    \begin{enumerate}
        \item Vrai.
              $\begin{aligned}[t]
                      10 & \stackrel{?}{=} 3\cdot 2 + 4 \\
                      10 & \stackrel{?}{=} 6 + 4        \\
                      10 & \ne 10
                  \end{aligned}$
        \item Vrai. Le graphe de cette fonction passe par $(0;0)$ car $f(x) = 2\cdot 0 = 0$.
        \item Faux. L'ordonnée à l'origine d'une fonction est l'ordonnée du point d'intersection de la fonction avec l'axe des ordonnées.
        \item Faux. La racine est la valeur de $x$ qui annule $f(x)$. Or $-2\cdot 3 - 6 = 12 \ne 0$.
    \end{enumerate}
